\section{Background and Related Work}

\subsection{Neural audio codecs}

Vector-quantised autoencoders~\cite{oord2017vqvae} turned audio into discrete tokens, and neural codecs made the tokens efficient: SoundStream~\cite{zeghidour2021soundstream}, EnCodec~\cite{defossez2022encodec}, HiFi-Codec~\cite{yang2023hificodec} and the Descript Audio Codec~\cite{kumar2023dac} compress waveforms with convolutional stacks descended from WaveNet~\cite{oord2016wavenet}, trained with reconstruction and adversarial losses. MiniMax Music 3's DAV follows the Descript encoder layout with the Snake activation~\cite{ziyin2020snake}, but it is a continuous variational autoencoder: its 128-channel posterior mean is the input of the diffusion transformer, and the eight discrete code streams are produced by a separate tokenizer that was not released.

\subsection{Residual vector quantisation}

Residual vector quantisation (RVQ) encodes a vector as a sum of codebook entries chosen in sequence, each codebook quantising the residual left by the previous ones~\cite{zeghidour2021soundstream}. The first codebooks carry coarse content and the later ones refine it, and the meaning of a code at depth $k$ depends on the codes chosen below it. Codec language models exploit this hierarchy with coarse-to-fine stages~\cite{borsos2023audiolm}, delayed patterns~\cite{copet2023musicgen} or a depth transformer~\cite{defossez2024moshi}. Our encoder does not learn a quantiser; it is a classifier over frozen codebooks whose entries we never observe, and Sec.~\ref{sec:analysis} shows that the conditional dependence across depth survives that framing.

\subsection{Audio tokenization for music generation}

Jukebox~\cite{dhariwal2020jukebox} modelled music with transformers over VQ-VAE codes. AudioLM~\cite{borsos2023audiolm} and MusicLM~\cite{agostinelli2023musiclm} separated semantic tokens, distilled from self-supervised speech representations such as HuBERT~\cite{hsu2021hubert}, from acoustic codec tokens. SpeechTokenizer~\cite{zhang2024speechtokenizer} and Mimi~\cite{defossez2024moshi} fold the semantic content into the first RVQ layer of the codec itself. MiniMax Music 3 uses one semantic codebook of 16,384 entries and seven acoustic codebooks of 1,024 entries, the same split, and renders the codes with latent diffusion~\cite{ho2020ddpm,lipman2023flow} over a transformer backbone~\cite{peebles2023dit,esser2024sd3}, as recent long-form music generators do~\cite{evans2024longform}.

\subsection{Existing music generation systems and model extraction}

Open music generators release their tokenizer with their language model~\cite{copet2023musicgen,agostinelli2023musiclm}; MiniMax Music 3~\cite{minimax2026music3} releases the generator but not the audio-to-code path. Recovering a model or one of its components from query access has been studied as model stealing: prediction APIs leak decision boundaries~\cite{tramer2016stealing}, black-box functionality can be copied into a knockoff network~\cite{orekondy2019knockoff}, and hidden dimensions of a production language model can be read out of its logits~\cite{carlini2024stealing}. Our problem is the benign counterpart: the released system exposes its outputs and its internal code distributions, and we recover the one component the release withholds by training on those outputs, with the temperature-scaled distillation of Hinton et al.~\cite{hinton2015distill} on the teacher's own samples~\cite{kim2016sequence,beyer2022patient}. All encoders are trained under maximal update parametrisation~\cite{yang2022mup}, and the gap we measure between teacher-forced and greedy decoding is the exposure bias of sequence models~\cite{bengio2015scheduled,ranzato2016sequence} across codebook depth rather than time.
